\section{Related Work}
\label{sec:related_work}

Our work is situated at the intersection of model merging, parameterized task consolidation, test-time adaptation, and signal processing. In this section, we discuss the relevant literature and highlight how our frequency-domain approach represents a paradigm shift.

\subsection{Post-hoc Model Merging}
With the advent of high-performance pre-trained foundation models, post-hoc model merging has emerged as an attractive, zero-shot alternative to joint multi-task training. Early methods such as Model Soups~\cite{wortsman2022model} and Task Arithmetic~\cite{ilharco2022editing} showed that pre-trained models fine-tuned from a shared initialization can be linearly combined in weight space to fuse their capabilities. However, simple linear interpolation often suffers from severe task interference due to representation clashes. To address this, several advanced merging algorithms have been proposed. TIES-Merging~\cite{yadav2023ties} resolves interference by signing and merging task vectors, filtering out redundant parameters with small magnitudes. DARE~\cite{dare2023} randomly drops delta-weights and rescales the remaining ones to preserve the original task capabilities. ZipMerge~\cite{zipmerge} co-optimizes merging and magnitude pruning to decouple task interference. While these methods improve upon uniform interpolation, they rely on pre-defined heuristics and lack the flexibility of dynamic, data-driven optimization. Crucially, SpectralMerge is conceptually orthogonal to, and fully compatible with, these sign-and-magnitude-based sparsification heuristics. Rather than acting as a replacement, SpectralMerge can operate in synergy with them: sign consensus and magnitude pruning (such as via TIES-Merging or DARE) can be first applied to isolate core non-interfering task updates, and SpectralMerge is subsequently utilized to optimize and regularize the layer-wise scaling of these filtered parameters in the frequency domain.

\subsection{Parameterized Model Merging}
To overcome the limitations of static merging, parameterized model merging dynamically optimizes layer-wise or parameter-wise combining coefficients using a small set of data. AdaMerging~\cite{adamerging} optimizes layer-wise coefficients by minimizing the entropy of model predictions on test-time data streams. RegCalMerge~\cite{regcalmerge} regulates and calibrates model merging via offline class-capacity normalization, preventing calibration drift. Similarly, other works have explored quantization-aware merging under non-differentiable operators (e.g., Q-Merge~\cite{qmerge}). However, a major bottleneck of these approaches is that they optimize coefficients as independent variables across layers. This spatial, coordinate-dependent representation suffers from a high dimensionality relative to the available data, causing severe overfitting to small data streams.

\subsection{Spatial Smoothing and Polynomial Trajectories}
To regularize parameterized merging and prevent overfitting, PolyMerge~\cite{polymerge} proposed to parameterize the merging coefficients as low-degree continuous polynomials over normalized network depth. This spatial smoothing approach acts as a spatial low-pass filter, forcing adjacent layers to have similar coefficients. However, we identify several fundamental limitations in PolyMerge's polynomial parameterization. First, polynomial basis functions (such as the standard power series) are highly collinear, meaning that the corresponding Vandermonde optimization matrix is extremely ill-conditioned. This ill-conditioning slows down optimization convergence and makes the search landscape unstable. Second, low-degree polynomials are highly sensitive to boundary conditions and cannot capture subtle, localized variations across layers without increasing the polynomial degree $d$. Yet, increasing $d$ quickly reintroduces the high-frequency degrees of freedom, leading back to overfitting. 

In contrast, our proposed \textbf{SpectralMerge} transitions the coefficient representation entirely into the frequency domain using the Discrete Cosine Transform (DCT-II). The DCT basis functions are strictly orthogonal, guaranteeing perfect numerical conditioning and an exceptionally stable optimization landscape. Furthermore, our spectral decay regularization allows us to softly penalize high frequencies, maintaining a high-fidelity representation of the underlying parameter sensitivity without boundary instabilities or collinearity.

\subsection{Test-Time Adaptation vs. Offline Few-Shot Tuning}
Our work also relates to test-time adaptation (TTA), where model parameters or merging coefficients are updated on the fly using unlabeled test streams. Typical TTA methods such as Tent~\cite{wang2020tent} minimize prediction entropy to adapt the model. However, recent work (e.g., ``The 'No-Data' Strawman''~\cite{nodata}) showed that online TTA under stream noise (such as label shift and bursty task streams) is highly unstable and often degrades performance. They proposed Offline Few-Shot Validation Tuning (OFS-Tune), demonstrating that optimizing low-dimensional coefficient profiles offline on tiny, labeled validation sets ($M \in [5, 50]$) completely outperforms online TTA under stream noise. In this work, we evaluate SpectralMerge under both online TTA and OFS-Tune paradigms. We show that while both paradigms benefit from frequency-domain modeling, SpectralMerge under the OFS-Tune regime achieves the highest level of robustness and generalization, completely refuting the Overfitting-Optimizer Paradox.
